% CROSCIM -- formulation et notations (v2, notations unifiees)
\documentclass[10pt]{article}
\usepackage[a4paper,margin=13mm]{geometry}
\usepackage[utf8]{inputenc}
\usepackage[T1]{fontenc}
\usepackage{amsmath,amssymb,mathtools,bm}
\usepackage[dvipsnames]{xcolor}
\usepackage[most]{tcolorbox}
\usepackage{enumitem,booktabs,array,tabularx,multirow}
\usepackage{helvet}
\renewcommand{\familydefault}{\sfdefault}
\usepackage[hidelinks]{hyperref}
\pagestyle{plain}
\setlength{\parindent}{0pt}
\setlength{\parskip}{3pt}
\setlist{nosep,leftmargin=12pt}

\definecolor{cC}{HTML}{1F5AA6}\definecolor{cM}{HTML}{1F8A5B}\definecolor{cH}{HTML}{D4561C}
\definecolor{ink}{HTML}{1D2230}\definecolor{muted}{HTML}{5B6475}
\definecolor{cS}{HTML}{7A4FA3}\definecolor{cP}{HTML}{B07A12}\definecolor{cO}{HTML}{2A7FB8}
\newcommand{\Cc}{{\color{cC}\mathcal{C}}}
\newcommand{\Mm}{{\color{cM}\mathcal{M}}}
\newcommand{\Hh}{{\color{cH}\mathcal{H}}}
\newcommand{\XC}{{\color{cC}X_{\mathcal{C}}}}
\newcommand{\XM}{{\color{cM}X_{\mathcal{M}}}}
\newcommand{\XH}{{\color{cH}X_{\mathcal{H}}}}
\newcommand{\dXM}{{\color{cM}\delta X_{\mathcal{M}}}}
\newcommand{\dXH}{{\color{cH}\delta X_{\mathcal{H}}}}
\newcommand{\SC}{{\color{cS}S_{\mathcal{C}}}}
\newcommand{\SM}{{\color{cS}S_{\mathcal{M}}}}
\newcommand{\psiC}{{\color{cP}\psi_{\mathcal{C}}}}
\newcommand{\psiM}{{\color{cP}\psi_{\mathcal{M}}}}
\newcommand{\PhiC}{\Phi_{\mathcal{C}}}
\newcommand{\PhiM}{\Phi_{\mathcal{M}}}
\newcommand{\ZC}{Z_{\mathcal{C}}}
\newcommand{\ZM}{Z_{\mathcal{M}}}
\newcommand{\hh}[1]{{\color{cO}h_{#1}}}
\newcommand{\Sop}{\mathcal{S}}
\newcommand{\Pup}{\Pi^{\uparrow}}
\newcommand{\Pdn}{\Pi^{\downarrow}}
\newcommand{\ch}{\mathrm{ch}}
\newcommand{\pa}{\mathrm{par}}
\newcommand{\nul}{\varnothing}
\newcommand{\E}{\mathbb{E}}
\newcommand{\obs}[1]{\widehat{#1}}

\newtcolorbox{bsec}[2][]{enhanced, breakable, colback=white, colframe=#2!70!white, boxrule=0.5pt, arc=2pt,
  left=6pt, right=6pt, top=8pt, bottom=4pt, fonttitle=\bfseries, coltitle=ink,
  colbacktitle=#2!10!white, title={#1}, attach boxed title to top left={xshift=5pt,yshift=-2.5pt},
  boxed title style={boxrule=0pt, colframe=white, arc=1pt}, before skip=9pt, after skip=4pt}
\newtcolorbox{note}{colback=black!3, colframe=black!3, boxrule=0pt, arc=1pt, left=4pt, right=4pt, top=2pt, bottom=2pt}

\begin{document}
\small\color{ink}
{\Large\textbf{CROSCIM\;---\;formulation et notations}}\\[2pt]
{\color{muted}Cascade de générateurs flow matching conditionnels sur une grille polaire stéréographique imbriquée, conditionnés par des contextes construits à la volée à partir d'observations hétérogènes et du forçage atmosphérique. Couleurs : niveaux {\color{cC}$\mathcal{C}$}, {\color{cM}$\mathcal{M}$}, {\color{cH}$\mathcal{H}$}\,; paramètres sous-maille explicites {\color{cS}$S$}\,; résidu latent {\color{cP}$\psi$}\,; information observationnelle {\color{cO}$h$}.}

% ======================================================================
\begin{bsec}[0 \;\; Glossaire (une notation = un objet)]{ink}
\renewcommand{\arraystretch}{1.12}
\begin{tabularx}{\linewidth}{@{}l X l@{}}
\toprule
\textbf{Symbole} & \textbf{Objet} & \textbf{Statut}\\\midrule
$\ell\in\{\Cc,\Mm,\Hh\}$, $\Omega_\ell$, $\Delta_\ell$ & niveaux, grilles, pas (32 / 4 / 1 km) ; $\ch(k)$ enfants, $\pa(i)$ parent & fixe\\
$R$, $P$, $\Pi_0=I-PR$ & restriction (moyenne pondérée océan), prolongation (injection), projecteur « détail » ; $RP=I$ & fixe\\
$\mathcal{V}_\ell$ & ensemble des variables d'état portées par le niveau $\ell$ ; $\mathcal{V}_\Hh\subseteq\mathcal{V}_\Mm\subseteq\mathcal{V}_\Cc$ & choix\\
$X=(\XC,\dXM,\dXH)$ & état : grande échelle + détails ($R\,\delta X=0$) & \textbf{généré}\\
$\SC,\ \SM$ & paramètres sous-maille \emph{explicites} (fractions, histogramme, anisotropie\dots) & \textbf{généré}, supervisé\\
$\Sop_\ell$ & opérateur statistique fixe : champ fin $\mapsto$ statistiques par maille de $\Omega_\ell$ & fixe\\
$\psiC,\ \psiM$ & résidu sous-maille \emph{latent} (petite dimension $d_\psi$) & appris\\
$\Phi_\ell=(S_\ell,\psi_\ell)$ & paramétrisation sous-maille du niveau $\ell$ & --\\
$Z_\Cc=(\XC,\SC[,\psiC])$, $Z_\Mm=(\dXM,\SM)$ & état augmenté par niveau (cible des générateurs) & \textbf{généré}\\
$Y=(Y_\Cc,Y_\Mm,Y_\Hh)$ & observations, rangées par résolution native ; $Y^{\rm in}\cup Y^{\rm tgt}$ : partition entrée / cible & données\\
$F$, $f$ & forçage ERA5 ($T_{2m},q_{2m},u_{10},v_{10}$ ; v1 : $T_{2m},\mathrm{MSL},u_{10},v_{10}$) ; $f=E_F(F)$ sur $\Omega_\Cc$ & données / appris\\
$E_\ell$, $a_\ell$ & encodeur local du niveau $\ell$ et sa sortie (grille $d\times\Omega_\ell$ + masque) & appris\\
$U_\ell$, $\hh{\ell}$ & bloc de fusion de la pyramide montante, information observationnelle du niveau $\ell$ & appris\\
$\Pup$, $\Pdn$ & montée (cross-attention sur les enfants), descente (injection + position) & appris\\
$\Psi_\ell$ & tête latente : $\psi_\ell=\Psi_\ell(\hh{\ell})$ & appris\\
$c_\ell$ & contexte de conditionnement du générateur $G_\ell$ & construit\\
$G_\ell$, $v_\theta^\ell$ & générateur FM du niveau $\ell$ et son champ de vitesse & appris\\
$\obs{X}_\ell$, $\obs{S}_\ell$, $m_\ell$ & vérités terrain construites \emph{à la volée} depuis $Y^{\rm tgt}$, masques de couverture & calculé\\
$\mathcal{N}_0$, $\mathcal{N}_1$ & réseau d'initialisation (Stage 1), réseau de transition (Stage 2) & appris\\
$\mu$, $w$ & disponibilité des observations (0/1), poids de rappel (guidance) & choix\\\bottomrule
\end{tabularx}
\end{bsec}

% ======================================================================
\begin{bsec}[1 \;\; Philosophie : un état multi-échelle et une chain rule grossier $\rightarrow$ fin]{ink}
\textbf{Grilles et opérateurs fixes.} $\Omega_\Cc\subset\Omega_\Mm\subset\Omega_\Hh$ imbriquées ($\times8$ puis $\times4$), arêtes alignées :
\[
(Rx)(k)=\frac{\sum_{i\in\ch(k)}\omega_i\,x(i)}{\sum_{i\in\ch(k)}\omega_i},\qquad P x=Sx+P_0(x-RSx),\qquad RP=I,\qquad \Pi_0=I-PR .
\]
$P_0$ : injection $(P_0x)(i)=x(\pa(i))$ ; $S$ : interpolation bilinéaire des parents (mailles terrestres ignorées). $RP=I$ reste exact puisque $RP_0=I$, mais $P\XC$ n'a plus de marches aux bords des mailles parentes : avec l'injection seule, les détails devaient effacer ces marches de 32 km et 4 km, et les champs générés restaient en blocs là où les détails ne sont pas supervisés.
\textbf{État.} Chaque niveau ne porte que ses variables $\mathcal{V}_\ell$ (section 2). Pour $v\in\mathcal{V}_\ell$ :
\[
\dXM=\XM-P\,\XC,\qquad \dXH=\XH-P\,\XM,\qquad R\,\delta X=0,\qquad
\XH^{v}=P\bigl(P\,\XC^{v}+\dXM^{v}\bigr)+\dXH^{v}\quad(v\in\mathcal{V}_\Hh).
\]
Une variable de $\mathcal{V}_\Mm\setminus\mathcal{V}_\Hh$ s'arrête à $\Mm$ ; de $\mathcal{V}_\Cc\setminus\mathcal{V}_\Mm$, à $\Cc$. Aucun décodeur appris : la recomposition est exacte.

\textbf{Paramètres sous-maille.} Ce que la grille d'un niveau ne résout pas est résumé par $\Phi_\ell=(S_\ell,\psi_\ell)$ :
\[
\SC=\Sop_\Cc(\text{champ fin dans la maille}),\qquad \SM=\Sop_\Mm(\cdot),\qquad \psi_\ell\in\mathbb{R}^{d_\psi}\ \text{(résidu latent)}.
\]
$\Sop_\ell$ est fixe, comme $R$, mais non linéaire. Les composantes \emph{additives} de $S$ (fractions, histogramme) s'emboîtent : $S_\Cc^{\rm frac}=R\,S_\Mm^{\rm frac}$.

\textbf{Loi cible et chain rule.} Avec les états augmentés $\ZC=(\XC,\SC)$, $\ZM=(\dXM,\SM)$ :
\[
p(Z\mid Y,F)=\underbrace{p(\ZC\mid Y,F)}_{G_\Cc}\;\cdot\;\underbrace{p(\ZM\mid \ZC,\,Y,F)}_{G_\Mm}\;\cdot\;\underbrace{p(\dXH\mid \ZC,\ZM,\,Y,F)}_{G_\Hh}.
\]
C'est une identité (aucune hypothèse). Le modèle remplace $(Y,F)$ par un contexte de taille fixe $c_\ell$ (section 3) :
\[
p_\theta(Z\mid Y,F)=p_\theta(\ZC\mid c_\Cc)\;p_\theta(\ZM\mid \ZC,\,c_\Mm)\;p_\theta(\dXH\mid \ZC,\ZM,\,c_\Hh).
\]
\begin{note}\footnotesize
\textbf{Lecture.} \emph{Grand $\rightarrow$ petit} : chaque niveau est conditionné par l'état déjà généré des niveaux plus grossiers (et par $\Pdn\Phi$, $\Pdn f$). \emph{Petit $\rightarrow$ grand} : deux canaux seulement --- (i) l'information observationnelle fine remontée par $\Pup$ dans $\hh{\Cc}$, (ii) les paramètres $\SC,\psiC$, qui font partie de l'état de $\Cc$. \textbf{Suffisance} : si $c_\ell$ contient toute l'information de $(Y,F)$ utile à $Z_\ell$ sachant $Z_{<\ell}$, la factorisation est exacte. Asymptotiquement, l'excès de loss FM d'un contexte $c$ par rapport aux données brutes vaut $\E\|v^\star_{Y}-v^\star_{c}\|^2\ge0$, nul ssi les lois conditionnelles coïncident : la loss mesure l'information perdue.
\end{note}
\end{bsec}

% ======================================================================
\begin{bsec}[2 \;\; Variables d'état et paramètres par niveau]{ink}
\renewcommand{\arraystretch}{1.15}\footnotesize
\begin{tabularx}{\linewidth}{@{}>{\bfseries}l >{\raggedright}p{3.1cm} >{\raggedright}X >{\raggedright\arraybackslash}p{4.1cm}@{}}
\toprule
& \textbf{état} $\mathcal{V}_\ell$ & \textbf{paramètres} $\Phi_\ell=(S_\ell,\psi_\ell)$ & \textbf{observations} $Y_\ell$\\\midrule
{\color{cC}$\mathcal{C}$ 32 km} & SIC$^\bullet$, $u^\bullet$, $v^\bullet$, SIT$^{\circ\!\bullet}$, neige$^{\circ}$ \newline {\color{muted}v1 : SIC, SIT, neige} &
{\color{cS}$\SC$} : eau libre$^{L}$, histogramme SIC$^{L}$, fraction de leads$^{L}$, densité de leads$^{F}$, anisotropie$^{F}$, déformation$^{F,\circ\bullet}$, glace mince$^{L,\circ\bullet}$, variance SIT / fraction déformée$^{F,\circ\bullet}$ ;\; {\color{cP}$\psiC$} &
SIC PMW (AMSR2, CIMR), dérive PMW, bouées, CS2SMOS / SMOS, neige \newline {\color{muted}v1 : CIMR (SIC, SIT mince), CRISTAL (neige)}\\[2pt]
{\color{cM}$\mathcal{M}$ 4 km} & $\delta$SIC$^\bullet$, $\delta u^{\circ\bullet}$, $\delta v^{\circ\bullet}$, $\delta$SIT$^{\circ\bullet}$ \newline {\color{muted}v1 : $\delta$SIC, $\delta$SIT} &
[{\color{cS}$\SM$} : eau libre, leads, variance SIC] (optionnel, depuis ASIP 500 m) &
dérive SAR, traces altimétriques (CS2, IS2), SIC SAR agrégée \newline {\color{muted}v1 : CRISTAL (SIT, traité en L2), ASIP agrégé}\\[2pt]
{\color{cH}$\mathcal{H}$ 1 km} & $\delta$SIC$^\bullet$ & --- (niveau le plus fin) & ASIP (SAR), dans les fauchées\\\bottomrule
\end{tabularx}\\[3pt]
{\footnotesize\color{muted}$^\bullet$ observé densément, $^{\circ\bullet}$ partiellement / clairsemé, $^{\circ}$ surtout inféré.\; $^{L}$ composante \emph{liée} à $\XC$ (la modifier modifie $\XC$ via $R$) ; $^{F}$ composante de \emph{forme} (modifiable à $\XC$ fixé). Le \emph{damage} n'est pas observé : il est remplacé par l'« état de fracturation » (densité de leads, anisotropie, déformation, texture). Le forçage $F$ n'est pas dans l'état : il entre par le contexte.}
\end{bsec}

% ======================================================================
\begin{bsec}[3 \;\; Le contexte $c_\ell$ : ce qu'il contient et comment on l'estime]{cO}
Le contexte d'un niveau a trois parts, de nature différente :
\[
c_\Cc=\bigl(\,\hh{\Cc},\ f\,\bigr),\qquad
c_\Mm=\bigl(\,\hh{\Mm},\ \Pdn f,\ \Pdn\PhiC\,\bigr),\qquad
c_\Hh=\bigl(\,\hh{\Hh},\ \Pdn f,\ \Pdn\PhiM\ \text{(ou }\Pdn\PhiC)\,\bigr).
\]
\begin{itemize}
\item {\color{cO}$\hh{\ell}$} : \textbf{information observationnelle} (latente, non interprétée). Disparaît quand il n'y a pas d'observation ($\hh{\ell}=\nul$). Rôle : rappel aux observations.
\item $\Phi_\ell=(S_\ell,\psi_\ell)$ : \textbf{paramétrisation sous-maille}. {\color{cS}$S_\ell$} : part \emph{physique explicite}, supervisée par des statistiques calculées à la volée sur les observations ; {\color{cP}$\psi_\ell$} : part \emph{latente inconnue}, estimée pour le conditionnement. À $\Cc$, $\SC$ est généré avec $\XC$ ; aux niveaux fins, $\Phi$ arrive d'en haut par $\Pdn$.
\item $f$ : \textbf{forçage}, connu, encodé une fois sur $\Omega_\Cc$ puis descendu.
\end{itemize}

\textbf{(a) Encodeurs locaux $a_\ell=E_\ell(Y_\ell^{\rm in})$ --- selon le type de donnée.} Chaque observation $n$ devient un token $e_n=\mathrm{MLP}\bigl[y_n,\ \sigma_n,\ \text{type}_n,\ x_n-x_{k(n)},\ t_n-t\bigr]$ (valeur, erreur, capteur, position dans la maille, décalage temporel).
\begin{center}\footnotesize\renewcommand{\arraystretch}{1.12}
\begin{tabular}{@{}lll@{}}\toprule
\textbf{type de données} & \textbf{exemples} & \textbf{encodeur}\\\midrule
points / traces / L2 irréguliers & altimétrie, bouées, dérive SAR & cross-attention par maille, requêtes apprises\\
champs denses déjà grillés & ASIP, SIC PMW grillée & CNN à pas (patch-embed masqué)\\
forçage régulier & ERA5 & CNN + regrillage $\rightarrow f\in\mathbb{R}^{d\times\Omega_\Cc}$\\\bottomrule
\end{tabular}
\end{center}
Pour les tokens irréguliers (softmax segmentée par maille $k$ ; maille vide $\Rightarrow$ token $a_\nul$ et masque) :
\[
a_\ell(k)=\sum_{n:\,x_n\in k}\operatorname*{softmax}_{n\in k}\!\Bigl(\tfrac{q^{\top}W_K e_n}{\sqrt d}\Bigr)\,W_V e_n .
\]
\textbf{(b) Pyramide montante (petit $\rightarrow$ grand).}
\[
\hh{\Hh}=a_\Hh,\qquad \hh{\Mm}=U_\Mm\bigl(a_\Mm,\ \Pup\hh{\Hh}\bigr),\qquad \hh{\Cc}=U_\Cc\bigl(a_\Cc,\ \Pup\hh{\Mm},\ f\bigr),
\]
\[
(\Pup h)(k)=W_O\Bigl[\sum_{i\in\ch(k)}\operatorname*{softmax}_{i\in\ch(k)}\Bigl(\tfrac{q_j^{\top}W_K\,(h(i)+p_i)}{\sqrt d}\Bigr)W_V\,(h(i)+p_i)\Bigr]_{j=1..K},\qquad (\Pdn g)(i)=W\,g(\pa(i))+p_i .
\]
$p_i$ : position de l'enfant dans le parent ; $q_j$ : $K$ requêtes apprises (« $K$ questions posées aux enfants ») ; enfants masqués exclus. Pour ASIP 500 m, $E_\Hh$ et $\Pup$ sont fusionnés par tuile $\Cc$ (CNN à pas $\rightarrow$ $\Pup$), uniquement sur les tuiles non vides : tout est calculé à la volée.

\textbf{(c) Part physique explicite : vérité terrain à la volée.}
\[
\obs{S}_\Cc=\Sop_\Cc\bigl(\obs{X}_\Hh\bigr)\ \text{(ou }\Sop_\Cc(\obs X_\Mm)\text{ pour la dérive)},\qquad m^{S}_\Cc(k)=\mathbb{1}\bigl[\mathrm{cov}(k)\ge\gamma\bigr],\qquad
\sigma_S^2(k)=\sigma_0^2+\frac{\sigma_{\rm samp}^2}{n_{\rm eff}(k)} .
\]
$\Sop$ est appliqué à une résolution effective commune ; les mailles partiellement couvertes ont une incertitude plus grande ($n_{\rm eff}$) ou sont masquées. Définitions v1 de $\Sop_\Cc$, $\gamma$ et $\sigma_S$ : section 8.

\textbf{(d) Part latente.} $\psi_\ell=\Psi_\ell(\hh{\ell})\in\mathbb{R}^{d_\psi}$, $d_\psi$ petit (4--8), régularisé (section 6) et décorrélé de $S_\ell$ pour que l'interprétable reste dans $S$.
\end{bsec}

% ======================================================================
\begin{bsec}[4 \;\; Flow matching en cascade : trois générateurs]{ink}
Pour chaque niveau, cible $z_1$ (état augmenté), bruit $\varepsilon\sim\mathcal{N}(0,I)$ (projeté par $\Pi_0$ pour les détails), interpolation linéaire :
\[
z_s=s\,z_1+(1-s)\,\varepsilon,\qquad u=z_1-\varepsilon,\qquad
\hat z_1=z_s+(1-s)\,v_\theta^{\ell}(z_s,s\mid\cdot)\ \ \text{(prédiction en un pas)}.
\]
\begin{center}\renewcommand{\arraystretch}{1.25}
\begin{tabular}{@{}l l l@{}}\toprule
\textbf{générateur} & \textbf{cible} $z_1$ & \textbf{conditionnement}\\\midrule
$G_\Cc$ & $\ZC=(\XC,\ \SC)$ & $c_\Cc=(\hh{\Cc},f)$\ \ [+ $\psiC=\Psi_\Cc(\hh{\Cc})$]\\
$G_\Mm$ & $\ZM=(\dXM,\ \SM)$ & $\XC^{\eta},\ \Pdn\PhiC,\ \hh{\Mm},\ \Pdn f$\\
$G_\Hh$ & $\dXH$ & $\XC^{\eta},\ \dXM^{\eta},\ \Pdn\PhiM,\ \hh{\Hh},\ \Pdn f$\\\bottomrule
\end{tabular}
\end{center}
$X^{\eta}=X+\eta$ : conditionnement grossier légèrement bruité (robustesse à ses erreurs). \textbf{À l'entraînement}, le conditionnement grossier provient d'un échantillon du niveau supérieur (modèle EMA, sans gradient), pas d'une vérité dense, qui n'existe pas. Avec probabilité $p_\nul$, chaque groupe de $c_\ell$ (en particulier $\hh{\ell}$, $\Phi$) est remplacé par $\nul$ : les deux têtes $v(\cdot\mid\dots,h)$ et $v(\cdot\mid\dots,\nul)$ sont apprises ensemble (guidance, section 7).

\textbf{Vérités terrain à la volée} (à chaque batch, depuis $Y^{\rm tgt}$ seulement) :
\[
\obs{X}_\ell=\mathcal{G}_\ell\bigl(Y_\ell^{\rm tgt}\bigr)\ \ \text{(moyenne pondérée par } \sigma_n^{-2}\text{ dans chaque maille)},\qquad
\obs{\delta X}_\ell=\obs{X}_\ell-P\,R\,\obs{X}_\ell\ \ \text{sur les parents couverts}\ (\ge\gamma).
\]
Le détail cible est défini \emph{par rapport à la moyenne des observations elles-mêmes} : il ne dépend pas de l'échantillon grossier et vérifie $R\,\obs{\delta X}=0$. Les masques $m_\ell^{v}$ (par variable $v$, niveau, pixel) valent 1 là où la cible existe.

\textbf{Échantillonnage.} Heun sur $s\in[0,1]$, dans l'ordre $\Cc\rightarrow\Mm\rightarrow\Hh$ ; $N$ tirages de $\varepsilon$ $\Rightarrow$ ensemble.
\end{bsec}

% ======================================================================
\begin{bsec}[5 \;\; Deux réseaux appris successivement]{cH}
\textbf{$\mathcal{N}_0$ --- initialisation (analyse).} Paramètres $\theta_0=\{E_\ell,U_\ell,\Pup,\Pdn,E_F,\Psi_\ell,G^0_\Cc,G^0_\Mm,G^0_\Hh\}$. Sur une fenêtre $[t-W,t]$ :
\[
Z^{t}\sim p_{\theta_0}\bigl(Z^{t}\mid Y^{\rm in}_{t-W:t},\,F_{t-W:t}\bigr)=p_{\theta_0}(\ZC^t\mid c_\Cc)\,p_{\theta_0}(\ZM^t\mid\ZC^t,c_\Mm)\,p_{\theta_0}(\dXH^t\mid\ZC^t,\ZM^t,c_\Hh).
\]
\textbf{$\mathcal{N}_1$ --- transition (évolution temporelle).} Initialisé depuis $\theta_0$ (encodeurs d'observations partagés, d'abord gelés). L'état précédent entre comme un flux d'observation supplémentaire, encodé par $E^{s}_\ell$ :
\[
Z^{t+\Delta}\sim p_{\theta_1}\bigl(Z^{t+\Delta}\mid Z^{t},\ \mu\,Y_{(t,t+\Delta]},\ F_{t:t+\Delta}\bigr),\qquad
a^{t+\Delta}_\ell=\bigl[E^{s}_\ell(Z^t_\ell),\ \mu\,E_\ell(Y_\ell)\bigr],\qquad \ZC=(\XC,\SC,\psiC).
\]
\begin{itemize}
\item \textbf{État d'entrée} $Z^t$ : un échantillon de $\mathcal{N}_0$ (l'analyse), puisqu'aucun état vrai dense n'existe.
\item \textbf{Cibles} à $t+\Delta$ : $\obs X$, $\obs S$ (section 4) et, pour le latent, $\obs\psi^{\,t+\Delta}=\bar\Psi\bigl(\hh{\Cc}^{t+\Delta}\bigr)$ (encodeur EMA, sans gradient), là où il y a des observations.
\item \textbf{$\mu\sim$ Bernoulli} : apprend à prévoir avec et sans observations nouvelles ; \textbf{pushforward} sur 2--4 pas (l'entrée du pas $k$ est la sortie du pas $k-1$) pour la stabilité du rollout.
\item \textbf{N1b} (options \texttt{increment}, \texttt{state\_input}, \texttt{raw\_obs}) : $G^1_\Cc$ génère l'incrément $Z^{t+\Delta}-Z^{t}$ (normalisé, $\psi$ compris), censuré pour que $Z^{t+\Delta}$ reste dans les bornes physiques : ne rien faire est la persistance ; l'état précédent entre aussi tel quel (pas seulement via $E^s$) et CIMR du jour non compressé (comme \texttt{raw\_days} pour $G^0_\Cc$). Le $\mathcal N_1$ de la v6, qui générait l'état complet à partir de $E^s(Z^t)$ et de $h_\Cc$ seul, était pire que la persistance dès $t+1$ en SIC (0{,}26 contre 0{,}19), avec ou sans observations ($w=1\approx w=0$) : il ne tirait parti ni de l'état ni des observations.
\item \textbf{N1c} (option \texttt{bg\_sigma}, terme d'ébauche) : là où une composante n'a pas de cible à $t+k$, sa cible est son état précédent, avec le plancher d'échelle $\sigma_{\rm bg}=0{,}5$ (unités normalisées) --- l'ébauche de la 4D-Var : sans observation l'état reste près de la persistance, à une dérive bornée près. Dans les rollouts N1b, $\SC$ (cibles seulement sous les passages ASIP retenus) était retiré au hasard à chaque pas (dispersion $\approx$ amplitude dès $t+1$) et, avec l'incrément, s'accumulait en marche aléatoire (\texttt{h\_thin} 1,0 $\to$ 2,5, \texttt{aniso} 1,1 $\to$ 3,1 en 7 jours) : le réseau recevait un état hors distribution et produisait du bruit. Les états d'entrée de l'entraînement sont aussi moins bruités (analyse à 20 pas de Heun, pushforward à 8).
\item \textbf{N1f} (options \texttt{inc\_scale}, \texttt{bg\_advect}) : (i) l'incrément est généré en unités de $0{,}2\times$ l'échelle de l'état --- avec l'échelle de l'état (SIC : identité) le bruit de départ du flow valait $\sim$10 fois les incréments journaliers, et son résidu, coupé d'un seul côté aux bornes (SIC 0 et 1), faisait baisser la banquise et monter l'eau libre à chaque pas : le lissage, présent même sur les mois d'entraînement (rapport prévision / persistance 1,15--1,8 en mars) alors que la neige, loin de sa borne, battait la persistance ; (ii) N1e : l'incrément et l'ébauche sont relatifs à l'état de la veille transporté par la dérive libre du jour, $\alpha\,\mathcal R(-\theta)\,\mathbf u_{10}$ ($\alpha,\theta$ appris, transport bilinéaire différentiable à $\Cc$), et cet état transporté entre aussi dans le conditionnement (étape « $\mathcal N_1$ avec $\mathcal T_{\mathbf u_0}$ » de la section 9). Les cibles faibles CIMR sont gardées (SIT mince).
\item \textbf{N1g} (option \texttt{warp\_mode}, $\sigma_{\rm bg}=0{,}05$) : rollouts du 24/12/2022 à $w=0$, N1c contre N1f (époque 12). (i) La SIC dérive vers la moyenne dans les deux (RMS normalisée 0{,}71 $\to$ 0{,}48 et 0{,}54 à 7 jours), dispersion SIC $\approx 0{,}1$ quel que soit le jour, prévision pire que la persistance là où celle-ci est presque parfaite (décembre, échéance 4 : 0{,}17 contre 0{,}05) : avec $\sigma_{\rm bg}=0{,}15$ (0{,}75 unité d'incrément) l'incrément était presque libre hors des passages ASIP ; $\sigma_{\rm bg}$ passe à 0{,}05. (ii) Le transport bilinéaire ($\sim$ une demi-maille $\Cc$ par jour) diffuse : sans transport (N1c) le bruit des incréments s'accumule (\texttt{sic\_std} 1,0 $\to$ 2,1, \texttt{aniso} 1,1 $\to$ 2,8), avec (N1f) la texture s'efface (0,74 et 0,62) ; le transport passe en bicubique (pondéré par l'océan, ramené dans les bornes physiques). Les pas d'entraînement dont la perte ou les gradients ne sont pas finis sont ignorés (N1f est passé à nan à l'époque 21).
\item \textbf{N1d} : pas de cibles faibles CIMR dans $\mathcal N_1$ (option \texttt{weak\_targets}) --- CIMR, plus bas qu'ASIP, tirait la prévision libre vers lui pas après pas (la SIC de la tuile baissait à chaque pas) ; l'ébauche prend le relais hors ASIP / CRISTAL, les zéros d'eau libre restent ; $\sigma_{\rm bg}=0{,}15$ (avec 0{,}5, h\_thin, h\_mid, sic\_std et l'anisotropie dérivaient encore comme une marche aléatoire, $\approx\sqrt{k}\,\sigma_{\rm bg}$).
\item \textbf{Mode A} (long terme) : seul $G^1_\Cc$ avance $(\XC,\SC,\psiC)$ ; les niveaux fins sont générés à la demande par $G^0_\Mm,G^0_\Hh$ avec $\hh{}=\nul$, conditionnés par $\Pdn\Phi$ et $\Pdn f$. \textbf{Mode B} (validation, régional) : la cascade complète à chaque pas, puis $\SC\leftarrow\Sop_\Cc(\XH)$.
\end{itemize}
\end{bsec}

% ======================================================================
\begin{bsec}[6 \;\; Les losses]{ink}
Notation : $\mathcal{L}_{\rm NLL}(u,v;\sigma,m)=\E\bigl[\sum m\odot\bigl(\tfrac12\bigl\|\tfrac{u-v}{\sigma}\bigr\|^2+\log\sigma\bigr)\bigr]$, avec $\sigma$ appris (tête log-échelle) et plancher connu pour $S$.
\[
\begin{aligned}
\mathcal{L}_X^{\ell}&=\mathcal{L}_{\rm NLL}\bigl(u^{X},v_\theta^{\ell,X};\ \sigma_\ell,\ m^X_\ell\bigr) &&\text{état, cible }\obs X_\ell\text{ ou }\obs{\delta X}_\ell\\
\mathcal{L}_S^{\ell}&=\mathcal{L}_{\rm NLL}\bigl(u^{S},v_\theta^{\ell,S};\ \sigma_S,\ m^S_\ell\bigr) &&\text{paramètres explicites, cible }\obs S_\ell=\Sop_\ell(\obs X_{\rm fin})\\
\mathcal{L}_{\rm cyc}&=\E\bigl\|\,\Sop_\Cc\bigl(\hat x_1^{\Hh}\bigr)-\SC\,\bigr\|^2_{m} &&\text{cohérence descente/montée ($\Sop$ doux)}\\
\mathcal{L}_\psi&=\E\|\psi\|^2+\lambda_\perp\bigl\|\mathrm{Cov}(\psi,S)\bigr\|_F^2 &&\text{goulot, décorrélation de }S\\
\mathcal{L}_\psi^{\rm dyn}&=\mathcal{L}_{\rm NLL}\bigl(u^{\psi},v_\theta^{\Cc,\psi};\ \sigma_\psi,\ m^{\psi}\bigr),\ \ \text{cible }\obs\psi^{\,t+\Delta} &&\text{latent dans la transition ($\mathcal{N}_1$)}
\end{aligned}
\]
\begin{center}\renewcommand{\arraystretch}{1.2}\footnotesize
\begin{tabular}{@{}lll@{}}\toprule
\textbf{réseau} & \textbf{loss totale} & \textbf{vérité terrain}\\\midrule
$\mathcal{N}_0$ & $\sum_\ell\bigl(\mathcal{L}_X^{\ell}+\lambda_S\mathcal{L}_S^{\ell}\bigr)+\lambda_c\mathcal{L}_{\rm cyc}+\lambda_\psi\mathcal{L}_\psi$ & $Y^{\rm tgt}_{t}$ (et $Y^{\rm in}$ avec poids $\lambda_{\rm in}<1$)\\
$\mathcal{N}_1$ & $\sum_{k=1}^{K}\gamma^{k-1}\Bigl[\sum_\ell\bigl(\mathcal{L}_X^{\ell}+\lambda_S\mathcal{L}_S^{\ell}\bigr)+\lambda_\psi\mathcal{L}_\psi^{\rm dyn}\Bigr]_{t+k\Delta}$ & $Y_{t+k\Delta}$, $\bar\Psi(\hh{}^{t+k\Delta})$\\\bottomrule
\end{tabular}
\end{center}
{\footnotesize\color{muted}La partition $Y=Y^{\rm in}\cup Y^{\rm tgt}$ (fauchées, traces ou blocs retirés au hasard à chaque batch) empêche le modèle de recopier ses entrées ; les scores sont toujours rapportés sur $Y^{\rm tgt}$. Aucune vérité n'est précalculée : grillage $\mathcal{G}_\ell$, détails, statistiques $\Sop$ et masques sont recalculés à chaque batch.}
\end{bsec}

% ======================================================================
\begin{bsec}[7 \;\; Scénarisation avec ces notations]{cS}
\textbf{Rollout (mode A).} À chaque pas, avec rappel optionnel aux observations ($w=0$ : simulation libre, $w=1$ : assimilation) :
\[
\ZC^{t+\Delta}\sim G^1_\Cc\bigl(\cdot\mid \ZC^{t},\,f^{t:t+\Delta}\bigr),\qquad
v_w=v\bigl(\cdot\mid \ZC^t,f\bigr)+w(t)\bigl[v\bigl(\cdot\mid \ZC^t,f,\hh{\Cc}\bigr)-v\bigl(\cdot\mid \ZC^t,f\bigr)\bigr].
\]
Statistiques + forçage pilotent la simulation \emph{non conditionnelle} ; les observations ajoutent un terme de rappel. Détails à la demande : $\dXM\sim G^0_\Mm(\cdot\mid\XC,\Pdn\PhiC,\Pdn f,\nul)$, puis $\dXH$.

\textbf{Interventions} (toujours suivies de la projection $\Pi_S$ : histogramme de moyenne $\mathrm{SIC}_\Cc$, fractions dans $[0,1]$) :
\begin{center}\renewcommand{\arraystretch}{1.3}\footnotesize
\begin{tabular}{@{}l l l@{}}\toprule
\textbf{intervention} & \textbf{opération} & \textbf{question}\\\midrule
(1) condition initiale & $\SC^{0}\leftarrow\Pi_S\bigl(\SC^0+\delta S^{F}\bigr)$, puis libre ; $\psiC^0$ regénéré sachant $\SC^0$ & mémoire sous-maille, relaxation\\
(2) forçage & $F\rightarrow\tilde F$ (+$\Delta T$, vents tournés, scénario) ; $\tilde f=E_F(\tilde F)$ & réponse de $\XC$ \emph{et} de $\SC$\\
(3) prescription & $\mathrm{do}(\SC^{t}=s^\star(t))$ après chaque pas & effet causal de $\SC$ sur $\XC$\\
(4) rappel & $\SC\leftarrow\SC+\tfrac{\Delta t}{\tau}(s^\star-\SC)$, ou guidance vers $s^\star$ & pousser sans casser la dynamique\\
(5) descente seule & $\dXH\sim G^0_\Hh(\cdot\mid\XC,\dXM,\Pdn\tilde\Phi,\Pdn f,\nul)$ & réalisations fines d'un $\tilde\Phi$ donné\\\bottomrule
\end{tabular}
\end{center}
$\delta S^{F}$ : modification de \emph{forme} (à $\XC$ fixé) ; modifier une composante liée ($^{L}$) impose de modifier $\XC$ de façon cohérente.

\textbf{Mesurer les rétroactions dans ce cadre.}
\[
I^{\uparrow}=\mathcal{L}_X^{\Cc}\bigl(\SC^{\perp}\!\rightarrow\nul\bigr)-\mathcal{L}_X^{\Cc},\qquad
I^{\downarrow}=\mathcal{L}_S^{\Cc}\bigl(f\text{ ou }\XC\!\rightarrow\nul\bigr)-\mathcal{L}_S^{\Cc},\qquad
R_{X\leftarrow S}(\tau)=\E_\varepsilon\bigl[\XC^{t+\tau}\mid\mathrm{do}(\SC^t+\delta S)\bigr]-\E_\varepsilon\bigl[\XC^{t+\tau}\mid\mathrm{do}(\SC^t)\bigr],
\]
avec $\SC^\perp=\SC-\E[\SC\mid\XC]$ et le \emph{même} $\varepsilon$ dans les deux rollouts (bruit apparié). Garde-fous : pente de contrôlabilité $\partial\,\Sop(\XH)/\partial S^\star\approx1$, modifications dans la gamme observée, placebo (permutation de $\SC$), cohérence mode A / mode B.
\end{bsec}

% ======================================================================
\begin{bsec}[8 \;\; Choix v1 (implémentation)]{cS}
\textbf{Sources v1} (\texttt{/Odyssey/public/CROSCIM\_dataset}) : ASIP (SIC 500 m, $\sigma$ = \texttt{standard\_deviation\_sic}, plancher 5\,\%) ; CIMR 5 km (SIC, SIT mince) ; CRISTAL 5 km (SIT, neige), \emph{traité comme du L2} : chaque maille 5 km valide est un point ; ERA5 sur la même grille 5 km. Pas de source de dérive en v1 : $u,v$ (et $\delta u,\delta v$) sont reportés en v2, ainsi que IST (non retenu) et la fraction de glace mince (demande le \texttt{sod} d'ASIP).

\renewcommand{\arraystretch}{1.15}\footnotesize
\begin{tabularx}{\linewidth}{@{}l l X@{}}
\toprule
\textbf{niveau} & \textbf{état v1} & \textbf{vérité terrain $\obs X$ (à la volée, depuis $Y^{\rm tgt}$, moyenne pondérée par $\sigma^{-2}$)}\\\midrule
{\color{cC}$\Cc$} & SIC, SIT, neige & SIC : ASIP agrégé ($R\,\mathcal{G}_\Hh$, mailles couvertes à $\ge50\,\%$ de leur océan), ailleurs CIMR du jour $t$ en cible faible (plancher d'échelle 0{,}15, comme $\sigma_S$) ; SIT : CRISTAL, ailleurs SIT mince CIMR ($\le1$ m) en cible faible (plancher 0{,}15 m) ; neige : CRISTAL ; SIT et neige sont des moyennes par maille : cible 0 là où la cible SIC est de l'eau libre ($\le0{,}05$ ; planchers 0{,}05 et 0{,}02 m). Cibles faibles et zéros non comptés dans les scores\\
{\color{cM}$\Mm$} & $\delta$SIC, $\delta$SIT & SIC : ASIP agrégé ($\mathcal{G}_\Mm=R\,\mathcal{G}_\Hh$) ; SIT : CRISTAL, 0 là où ASIP voit de l'eau libre (plancher 0{,}05 m)\\
{\color{cH}$\Hh$} & $\delta$SIC & ASIP (blocs $2\times2$ exacts de la grille 500 m)\\\bottomrule
\end{tabularx}\\[2pt]
\textbf{Grillage $\mathcal{G}_\ell$ sans interpolation.} Champs alignés (ASIP $\to\Hh$) : moyenne par blocs ; champs denses non alignés (CIMR, ERA5 5 km $\to\Cc$) : remappage exact par recouvrement de surfaces (matrices 1D séparables $A_y\,X\,A_x^\top$) ; points (CRISTAL) : affectation à la maille contenant le point. Poids de $R$ : fraction d'océan $\omega_\Hh$ (Natural Earth 10 m, sur-échantillonnée $\times4$), sommée vers $\Mm$ et $\Cc$.

\textbf{Partition $Y^{\rm in}\cup Y^{\rm tgt}$} (jour $t$ seulement, les jours précédents sont des entrées) : ASIP par passage (temps d'acquisition chaînés par écarts $<10$ min) ; CRISTAL par blocs de $8\times8$ mailles $\Cc$ (256 km) ; chaque unité va dans $Y^{\rm tgt}$ avec probabilité 0{,}3. \textbf{Une seule cible par variable} : CIMR n'est jamais une cible, il entre en entier dans $Y^{\rm in}$ comme un \emph{prior} (esprit assimilation) ; le réseau apprend, là où ASIP est en $Y^{\rm tgt}$, la correction (non linéaire) de CIMR vers ASIP, et sans ASIP il s'appuie sur CIMR corrigé. La couverture des cibles $\Cc$ suit celle d'ASIP et de CRISTAL.

\textbf{Détails cibles.} $\obs{\delta X}=\obs X-PR\,\obs X$ avec $R$ pondéré par $\omega\cdot m$ (mailles observées), sur les parents dont la fraction d'océan observée est $\ge\gamma$ : $\gamma_{\rm SIC}=0{,}5$ (fauchées SAR larges), $\gamma_{\rm SIT}=0{,}05$ (une trace CRISTAL, points tous les 5 km, touche $\approx$ 6 des 64 mailles $\Mm$ d'une maille $\Cc$ ; avec $0{,}1$, 0,5\,\% des mailles seulement avaient une cible).

\textbf{$\SC$ v1 : 9 canaux}, calculés par $\Sop_\Cc$ sur $\obs X_\Hh^{\rm SIC}$ (résolution effective commune : 1 km) et sur $\obs X_\Mm^{\rm SIT}$ :
\begin{center}\footnotesize\renewcommand{\arraystretch}{1.2}
\begin{tabularx}{\linewidth}{@{}l l X l@{}}\toprule
\textbf{canal} & \textbf{type} & \textbf{définition (moyennes pondérées par $\omega$ sur les enfants observés)} & $\sigma_{\rm samp}(\obs S)$\\\midrule
$h_0..h_3$ & $L$ & histogramme SIC, classes $[0,.15)$, $[.15,.5)$, $[.5,.85)$, $[.85,1]$ ; $h_0$ = eau libre ; en banquise compacte, $h_0+h_1$ = fraction de leads & $\sqrt{h(1-h)}$\\
sic\_std & $F$ & écart-type de la SIC à 1 km dans la maille & $\obs S/\sqrt2$\\
edge\_density & $F$ & fraction des paires de voisins (dans la maille) avec $|\Delta\mathrm{SIC}|>0{,}1$ : densité de bords de leads / floes & $\sqrt{e(1-e)}$\\
aniso\_c, aniso\_s & $F$ & $\bigl(\bar J_{xx}-\bar J_{yy},\ 2\bar J_{xy}\bigr)/(\bar J_{xx}+\bar J_{yy}+e_0)$, tenseur de structure des gradients aux coins des blocs $2\times2$ (non biaisé pour un bruit non corrélé), $e_0=10^{-3}$ (niveau du bruit de restitution) : anisotropie $\times$ orientation ($2\theta$) & $1-a^2$\\
sit\_std & $F$ & écart-type de la SIT CRISTAL sur les mailles $\Mm$ observées ($\ge4$) & $\obs S/\sqrt2$\\\bottomrule
\end{tabularx}
\end{center}
Masque : $m^S=\mathbb 1[\mathrm{cov}\ge0{,}5]$ (canaux ASIP), $\mathbb 1[n_\Mm\ge4]$ (sit\_std) ; $n_{\rm eff}=\sum\omega\,m/4$ (décorrélation $\approx2$ km) ou $n_\Mm$ ; $\sigma_S^2=\sigma_0^2+\sigma_{\rm samp}(\obs S)^2/n_{\rm eff}$, $\sigma_{\rm samp}$ = erreur-type d'un échantillon indépendant (proportion, écart-type, coefficient de type corrélation, $a$ = module de l'anisotropie), $\sigma_0$ = 10\,\% de l'écart-type mesuré du canal (janvier--octobre 2022), 20\,\% pour les canaux de texture (sic\_std, edge\_density, $a$, sit\_std), dont la variabilité d'un jour à l'autre est en partie du bruit d'échantillonnage. Des constantes « pire cas » ($\sigma_{\rm samp}=0{,}5$) donnaient à edge\_density et sit\_std un plancher 2 à 4 fois plus grand que leur variabilité : ces canaux n'étaient pas appris.

\textbf{Échantillons et fenêtres.} Tuiles de $64\times64$ mailles $\Cc$ (2048 km : $512^2$ en $\Mm$, $2048^2$ en $\Hh$), à fraction d'océan $\ge0{,}2$ ; $\mathcal N_0$ : $W=15$ jours ; ASIP entre à $\Hh$ pour les 3 derniers jours et à $\Mm$ (une carte par jour, moyenne pondérée par $\sigma^{-2}$ et somme des poids) sur les 15 jours : le détail à 1 km des jours anciens est périmé par la dérive, leur couverture et leur état moyen restent utiles. Les cartes $\Mm$ des jours $<t$ sont stockées à la conversion (aucune partition ne s'applique à ces jours : l'agrégat est le même à chaque batch ; l'agréger dans les workers limitait le débit à 0,5 s par batch), celle du jour $t$ est calculée sur GPU après la partition ; $\mathcal N_1$ : $\Delta=1$ jour, rollout appris jusqu'à 7 jours, avec un curriculum de 1 à 7 pas.

\textbf{Stockage.} Conversion unique des NetCDF en Zarr chunké par jour (recadrage au domaine, valeurs natives, aucun regrillage). Les workers ne font que découper des tranches ; partition, $\mathcal{G}_\ell$, masques, $\obs{\delta X}$ et $\obs S$ sont calculés sur GPU après le transfert du batch.

\textbf{Implémentation v1 (choix non fixés par les sections 1--7).}
\begin{itemize}
\item \emph{Encodeurs} : CRISTAL en tokens (SIT à $\Mm$, neige à $\Cc$, softmax segmentée par maille, fenêtre complète avec $t_n-t$) ; CIMR, ASIP agrégé à $\Mm$ et ERA5 par CNN masqués ; ASIP à $\Hh$ fusionné avec $\Pup_{\Hh\to\Mm}$ (CNN à pas par maille $\Cc$ non vide) ; $\Pup_{\Mm\to\Cc}$ à $K=4$ requêtes ; $d=32$, $d_\psi=4$.
\item \emph{Prolongation lisse} $P=S+P_0(I-RS)$ (section 1), calculée sur la tuile entière (resp. le crop $\Mm$) avant de découper, pour que les crops aient leurs vrais voisins ; elle remplace l'injection depuis la deuxième passe 2022 (G$_\Mm$, G$_\Hh$, pyramide et $\mathcal N_1$ réentraînés ; G$_\Cc$ en est indépendant).
\item \emph{Troisième passe 2022} (G$_\Cc$ moins bon que la persistance CIMR, champs fins encore en blocs) : (i) les moyennes CIMR brutes des 3 derniers jours (et leurs masques) entrent telles quelles dans le conditionnement de G$_\Cc$, à côté de $\hh\Cc$ ; (ii) $\Pdn$ et le conditionnement neige de G$_\Hh$ passent de l'injection à l'interpolation bilinéaire ; (iii) à l'entraînement de G$_\Mm$, G$_\Hh$ et de la pyramide, 80\,\% des tuiles sont tirées proportionnellement à leur couverture ASIP du jour $t$ (20\,\% uniformes) ; (iv) nouveaux $\sigma_0$ de $\sigma_S$ et normalisation de $\SC$ ; tout est réentraîné depuis G$_\Cc$. \emph{Quatrième passe} : cibles uniques par variable (tableau ci-dessus) ; avant, la cible SIC à $\Cc$ était CIMR, en désaccord avec ASIP de 0{,}08 à 0{,}21 en RMSE (biais ASIP$-$CIMR de $+0{,}01$ à $+0{,}11$, décembre 2022), écart qui dominait les erreurs des champs absolus à $\Mm$ et $\Hh$. \emph{Cinquième passe} : (i) sans cible SIC loin du bord (où ASIP ne passe pas), G$_\Cc$ dérivait en eau libre (SIC 0{,}1--0{,}2) : CIMR y redevient une cible faible ; (ii) \emph{G$_\Mm$, G$_\Hh$ correcteurs} : avec $R\,\delta X=0$ et des cibles $Y-P'RY$, G$_\Mm$ ne voyait jamais l'erreur de l'échantillon $X_\Cc$ qui le conditionne, et toute erreur maille à maille de $X_\Cc$ restait en marches de 32 km à $\Mm$ et $\Hh$ (les détails ne battaient pas $\delta X=0$). Désormais $\delta X_\Mm=X_\Mm-P'X_\Cc^{\rm éch}$ et $\delta X_\Hh=X_\Hh-P'X_\Mm^{\rm éch}$, cibles $Y-P'X^{\rm éch}$, sans projection $\Pi_0$, censurées aux bornes physiques du champ corrigé ; $X_\Cc$ reste l'échantillon de G$_\Cc$ ($R\,X_\Mm$ en est la version cohérente) : la décomposition exacte de la section 1 devient un a priori corrigé par les niveaux fins, comme CIMR l'est par ASIP. Échelles des corrections : premières estimations (0{,}1 ; 0{,}08 m ; 0{,}05).
\item \emph{Générateurs fins sur crops} alignés sur les mailles $\Cc$ : $G_\Mm$ sur $128^2$ ($16\times16$ mailles $\Cc$), $G_\Hh$ sur $256^2$ ($8\times8$). Le bruit \emph{et la vitesse} sont projetés par $\Pi_0$ : $R\,\delta X=0$ exactement (jusqu'à la quatrième passe ; correcteurs ensuite, voir ci-dessus). Tuile entière = crops juxtaposés (raccords possibles aux bords des crops). Pas de bornes sur les $\delta$ : seuls les champs recomposés peuvent sortir légèrement de $[0,1]$.
\item \emph{$\psi_\Cc=\Psi_\Cc(\hh\Cc)$} est déterministe dans $\mathcal N_0$ (il conditionne $G_\Mm$, $G_\Hh$ via $\Pdn$), généré par $G^1_\Cc$ dans $\mathcal N_1$.
\item \emph{$\mathcal L_{\rm cyc}$} : $\Sop_\Cc$ doux (indicatrices sigmoïdes, $\tau=0{,}05$) appliqué à la prédiction en un pas de $\Hh$, comparé au $\SC^\eta$ qui l'a conditionnée, pondéré par $s$.
\item \emph{Dropout de groupe} ($p_\nul=0{,}1$) : $(\hh{},\Phi,f)$ dans $\mathcal N_0$ ; $(\hh{},f,X,S,\psi)$ dans $\mathcal N_1$, où $\hh{}=\nul$ aussi si $\mu=0$. $I^\uparrow$ est estimé en remplaçant tout le groupe $S$ par $\nul$ : le modèle marginalise alors $S$ sachant $X$, ce qui retire $S^\perp$.
\item \emph{$\Pi_S$} : bornes physiques et histogramme normalisé (histogramme vide $\Rightarrow$ classe de $\mathrm{SIC}_\Cc$). La cohérence entre la moyenne de l'histogramme et $\mathrm{SIC}_\Cc$ n'est pas imposée en v1.
\item \emph{$\lambda_{\rm in}$} (sixième passe, correcteurs) : l'ASIP du jour $t$ de $Y^{\rm in}$ est aussi cible de G$_\Mm$ et G$_\Hh$, avec le poids $\lambda_{\rm in}=0{,}3$ (1 pour $Y^{\rm tgt}$, moyenne pondérée par canal). Sans lui, la partition rend disjoints les pixels vus et les pixels cibles : aucune perte là où le générateur voit une observation, et l'évaluation sans retenue montrait qu'il ignorait l'ASIP fourni. Les scores restent calculés sur $Y^{\rm tgt}$ seul.
\item \emph{Texture de la SIT à $\Mm$} ($\lambda_{\rm cyc}^{\rm SIT}=0{,}1$, v7) : l'écart-type de la SIT générée en un pas dans chaque maille $\Cc$ doit égaler le sit\_std de $\SC$ qui conditionne G$_\Mm$ (comme $\mathcal L_{\rm cyc}$ pour la SIC). En v6 les membres avaient une variabilité sous-maille 2--3 fois celle de CRISTAL et 3--8 fois le sit\_std prédit par G$_\Cc$, et les corrections de SIT déplaçaient la moyenne des mailles de 0{,}2--0{,}4 m dans toute la glace (évaluation sur 12 dates, recadrages choisis par CRISTAL) : utile quand G$_\Cc$ est très faux, nuisible sinon. D'où aussi un \emph{a priori de décalage} ($\lambda_{\rm shift}^{\rm SIT}=0{,}1$, v7) : les moyennes par maille $\Cc$ de la correction de SIT prédite en un pas (en unités de son échelle) sont ramenées vers 0, pondérées par le pseudo-temps ; les cibles CRISTAL l'emportent là où elles exigent un décalage.
\item \emph{Non implémentés en v1} : $\SM$, régénération de $\psi^0$ après une intervention sur la condition initiale.
\end{itemize}
\end{bsec}

% ======================================================================
\begin{bsec}[9 \;\; Dérive apprise des seules observations (v2, proposition non implémentée)]{cO}
\textbf{Contrainte.} Le produit n'utilise que ASIP, CIMR, CRISTAL et le forçage ERA5 : aucune observation de dérive, donc aucune cible pour $\mathbf u=(u,v)$. La dérive est la \emph{variable de contrôle} qui explique au mieux la succession des observations (esprit 4D-Var : un modèle de transport, des observations à des instants différents). Elle n'a de sens physique que si la structure du modèle l'oblige à agir comme un transport ; sinon elle n'est qu'un latent de plus.

\textbf{Information disponible.} ASIP (SIC 500 m, heure d'acquisition par pixel) : déplacement des bords, chenaux et floes entre deux jours, avec l'intervalle exact $\Delta t(x)$ ; le L3 n'a qu'une valeur par pixel et par jour (pas de paires intra-journalières au même endroit). CIMR : déplacement des gradients de SIC à 5 km. CRISTAL : traces éparses, faible. ERA5 : a priori de dérive libre, $\mathbf u_0=\alpha\,\mathcal R(-\theta)\,\mathbf u_{10}$ ($\alpha\approx1$--$2\,\%$, $\theta\approx10$--$25^\circ$ vers la droite ; $\mathbf u_{10}$ selon les axes de la grille).

\textbf{Transport en forme flux} (moyennes par maille, section 8) :
\[
\partial_t A+\nabla\!\cdot(\mathbf u A)=S_A,\qquad \partial_t h+\nabla\!\cdot(\mathbf u h)=S_h,\qquad \partial_t h_s+\nabla\!\cdot(\mathbf u h_s)=S_{h_s},
\]
discrétisé en volumes finis (ou semi-lagrangien conservatif) à $\Mm$ : une dérive de 5--20 km/jour est sous-maille à $\Cc$. La divergence $\nabla\!\cdot\mathbf u$ porte l'ouverture des chenaux et la convergence ; $S$ porte la thermodynamique.

\textbf{Paramétrisation.} $\mathbf u=\mathbf u_0(\mathbf u_{10};\alpha,\theta)+\mathbf u_{\rm res}$, avec $\alpha,\theta$ appris (d'abord globaux, puis fonctions de $h$, $A$) et $\mathbf u_{\rm res}=U_\theta(\ZC,f)$ à $\Cc$, prolongé par $P'$ vers $\Mm$ (puis, plus tard, $\delta\mathbf u_\Mm$).

\textbf{Apprentissage dans $\mathcal N_1$} (correcteur du transport, comme G$_\Mm$, G$_\Hh$ corrigent $P'X$) :
\[
X^{t+\Delta}=\underbrace{\mathcal T_{\mathbf u,\Delta}\bigl(X^t_\Mm\bigr)}_{\text{prior transporté}}+\ \delta X,\qquad \delta X\sim G^1(\cdot\mid \mathcal T_{\mathbf u}X^t,\ \ZC^t,\ \mu Y,\ f),
\]
cibles et pushforward inchangés (section 5) ; $\mathbf u$ est appris de bout en bout par la perte sur les observations futures. Régularisations : $\lambda_u\|\mathbf u_{\rm res}\|^2$ (rappel vers la dérive libre), $\lambda_\nabla\|\nabla\mathbf u\|^2$ (problème d'ouverture : là où la glace est sans structure seule la composante normale aux gradients est contrainte), $\lambda_\delta\|\delta X\|^2$ (sinon la correction absorbe le déplacement et $\mathbf u$ n'apprend rien). Les termes $S$ sont conditionnés par t2m, $h$, $h_s$ pour séparer dynamique et thermodynamique.

\textbf{Diagnostics et validation sans données de dérive.} (0) Faisabilité, avant tout apprentissage (\texttt{scripts/drift\_feasibility.py}) : SIC ASIP de $t-1$ transportée par la dérive libre vs persistance, à $\Mm$, dans la zone du bord ; suivi par blocs de 64 km sur ASIP seul et comparaison à la dérive libre (rapport de vitesse, angle, corrélation vectorielle), ce qui dit aussi quel repère ont $u_{10},v_{10}$ dans les stores. (1) Passages ASIP retenus : $\mathcal T_{\mathbf u}$ doit battre la persistance et la dérive libre seule. (2) Cohérence physique : vitesse/vent plus faible pour la glace épaisse et compacte, plus forte au bord. (3) Carte d'identifiabilité ($|\nabla\mathrm{SIC}|$, couverture à deux jours). Une dérive externe (OSI SAF, bouées) ne servirait qu'à l'évaluation, hors du produit.

\textbf{Faisabilité 2022} (42 paires $t-1,t$, ASIP à $\Mm$). $u_{10},v_{10}$ des stores sont selon les axes de la grille (repère grille : compétence positive dans 95\,\% des paires, contre 50\,\% en est/nord). Dérive libre vs persistance : $-18\,\%$ d'erreur quadratique dans la zone du bord ($\alpha=2\,\%$), $-14\,\%$ partout ($\alpha=1\,\%$ : la banquise compacte freine). Suivi par blocs sur ASIP seul : 51\,\% des blocs texturés suivis, corrélation vectorielle 0{,}44 avec la dérive libre, rotation réelle $\approx10^\circ$ à droite du vent. La dérive est identifiable là où les fauchées de deux jours se recouvrent et où la glace a de la structure ; ailleurs, l'a priori du vent.

\textbf{Étapes.} (0) faisabilité (faite) ; (1) \emph{fenêtre de $\mathcal N_0$ transportée} (option \texttt{advect\_window}, pipeline v7) : chaque jour $k$ de la fenêtre (CIMR, ASIP à $\Mm$ et $\Hh$, points CRISTAL) est déplacé vers $t$ par $\sum_{j=k}^{t-1}\alpha\,\mathcal R(-\theta)\,\bar{\mathbf u}_{10,j}\,\Delta t$ (vitesse au point d'arrivée), échantillonnage bilinéaire différentiable, $\alpha,\theta$ appris par la perte de $\mathcal N_0$ (gelés avec G$_\Cc$ en phase 4), jour $t$ inchangé ; (2) $\mathcal N_1$ avec $\mathcal T_{\mathbf u_0}$ ; (3) $\mathbf u_{\rm res}$ à $\Cc$ régularisé ; (4) $\mathcal N_1$ à $\Mm$ pour SIC (et $\delta\mathbf u_\Mm$).
\end{bsec}

\end{document}
